%قبل از شروع مسئله ترجیحاً چیزی ننویسید. اما اگر خواستید بنویسید، در قالب لیستی (فقرات و شمارش) بنویسید.
%برای اضافه کردن نمره، بعدی از شروع‌مسئله، [x] را اضافه کنید و به جای x  نمره مورد نظر را قرار دهید. برای زیرمسئله‌ها هم به همین صورت است.
%\شروع‌مسئله[10]
\شروع‌مسئله[25]
در این سؤال روی محور $x$ها تعدادی بازه قرار دارد که هر بازه یک وزن دارد. در هر مرحله یا یک بازه اضافه می‌شود و یا یک نقطه روی محور داده می‌شود و برنامۀ شما باید مجموع وزن بازه‌هایی که این نقطه آن‌ها را قطع می‌کند را برگرداند. توجه کنید اگر نقطه سر یا ته بازه باشد آن را قطع نمی‌کند.

در هر خط ورودی یا 
\lr{\texttt{add a b w}} داده می‌شود که به معنی این است که بازۀ $[a, b]$ با وزن $w$ اضافه شود. همچنین 
\lr{\texttt{ask p}} نقطۀ $p$ را مورد پرسش قرار می‌دهد.
$a$، $b$، $w$ و $p$ همه اعداد صحیح مثبت و کمتر از $1,000,000$ هستند.

شما در این سؤال باید از داده‌ساختار درخت فنویک (درخت اندیس داده شدۀ دودویی) که در کلاس مطرح شد استفاده کنید.

\begin{center}
	\begin{tabular}{|c|c|}
		\hline
		خروجی                 
		& ورودی    \\ \hline
		\parbox[t]{5cm}{\normalsize\raggedright{\texttt{150\\100\\60\\160}}\vspace{0.2cm}} &
		\parbox[t]{5cm}{\normalsize\raggedright{\lr{\texttt{add 1 10 50\\add 5 20 100\\ask 6\\ask 10\\add 2 8 10\\ask 3\\ask 6}}}\vspace{0.2cm}}\\\hline
	\end{tabular}
\end{center}

%\زیرمسئله % مسئلۀ ۱پ
%\زیرمسئله % مسئلۀ ۱ت


\solution{
برای راحتی اندیس‌ها را در ۲ ضرب می‌کنیم و هر عمل را با استفاده از درخت فنویک به صورت زیر انجام می‌دهیم. عمل \متن‌لاتین{\متن‌تایپ{add a b w}} به دو عمل \متن‌لاتین{\متن‌تایپ{update(2a+1, w)}} و \متن‌لاتین{\متن‌تایپ{update(2b-1, -w)}} در درخت فنویک تبدیل می‌شود. عمل \متن‌لاتین{\متن‌تایپ{ask p}} نیز به عمل \متن‌لاتین{\متن‌تایپ{sum(2p)}} در درخت فنویک تبدیل می‌شود.

کد پایتون این سؤال از خانم مهتاب بیگوردی ضمیمه شده است.

}




%%%%%%%%%%%%%%%%%%%%%%%%%%%%%%%%%%%%%%%%%%%%%%%%%%%%%
% برخی دستورات مفید برای تمرین‌ها. %
%%%%%%%%%%%%%%%%%%%%%%%%%%%%%%%%%%%%%%%%%%%%%%%%%%%%%

% فرمول‌های نمونه:
%$f(x) = \Theta(x)$
%$T(x, y) \leq \log(x) + 2^y + \binom{2n}{n}$
% {\tt code\_function}


% لیست موارد بدون شماره‌گذاری:
%\شروع{فقرات}
%    \فقره مورد اول 
%        \شروع{فقرات}
%            \فقره مورد اول داخلی 
%                \شروع{فقرات}
%                    \فقره اولین مورد خیلی داخلی 
%                    \فقره دومین مورد خیلی داخلی 
%                \پایان{فقرات}
%            \فقره مورد دوم داخلی 
%        \پایان{فقرات}
%    \فقره مورد دوم 
%\پایان{فقرات}

% لیست موارد شماره‌گذاری شده:
%\شروع{شمارش}
%    \فقره یک مورد 
%    \فقره یک مورد دیگر 
%    \فقره باز هم یک مورد دیگر
%\پایان{شمارش}

% برای نوشتن یک فرمول چند خطی که یک شماره بخوره. برای حالت‌های دیگر به سایت زیر مراجعه کنید:
% https://www.overleaf.com/learn/latex/Aligning_equations_with_amsmath
%\begin{align} 
%    \begin{split}
%        (x+y)^3 &= (x+y)^2(x+y) \\
%                &= (x^2+2xy+y^2)(x+y) \\
%                &= (x^3+2x^2y+xy^2) + (x^2y+2xy^2+y^3) \\
%                &= x^3+3x^2y+3xy^2+y^3
%    \end{split}                                 
%\end{align}

% نوشتن ماتریس و آرایه دوبعدی
%\begin{align}
%    A = 
%    \begin{bmatrix}
%        A_{11} & A_{21} \\
%        A_{21} & A_{22}
%    \end{bmatrix}
%\end{align}

% اضافه کردن تصویر با شرح. برای اضافه کردن تصویر بدون شرح، فقط دستور درج‌تصویر را استفاده کنید.
%\شروع{شکل}[h!]
%\تنظیم‌ازوسط
%  \درج‌تصویر[پهنا=0.5\پهنای‌متن]{img.jpg}
%\شرح{شرح تصویر}
%\برچسب{تصویر:۱}
%\پایان{شکل}
%\پایان{زیرمسائل}

%کد زیر را ببینید.
%%\lstinputlisting{source_filename.py}
%\begin{latin}
%\begin{lstlisting}
%# Import the modules
%import sys
%import random
%
%ans = True
%
%while ans:
%question = raw_input("Ask the magic 8 ball a question: (press enter to quit) ")
%
%answers = random.randint(1,8)
%
%if question == "":
%sys.exit()
%
%elif answers == 1:
%print "It is certain"
%
%elif answers == 2:
%print "Outlook good"
%
%elif answers == 3:
%print "You may rely on it"
%
%elif answers == 4:
%print "Ask again later"
%
%elif answers == 5:
%print "Concentrate and ask again"
%
%elif answers == 6:
%print "Reply hazy, try again"
%
%elif answers == 7:
%print "My reply is no"
%
%elif answers == 8:
%print "My sources say no"
%\end{lstlisting}
%\end{latin}
